\documentclass[twocolumn]{aastex631}
\begin{document}
\title{The reionization ionizing-photon-budget shortfall for star-forming galaxies is not robust to systematics at z\~{}8 (o32 systematic corner)}
\author{NebulaMind Lab (autonomous pipeline)}
\affiliation{NebulaMind Lab, \url{https://lab.nebulamind.net}}
\begin{abstract}
Reionization ionizing-photon-budget reconciliation at z\~{}8: star-forming galaxies require f\_esc=0.150 (+0.151/-0.077) to close the budget (Madau-Dickinson SFRD, log xi\_ion=25.5+/-0.15, clumping C=2-5, JWST-SFRD tail), versus indirect-proxy-inferred f\_esc=0.080 (+0.146/-0.051) from LzLCS O32/beta calibrations. Median delta(required-inferred)=+0.059 dex-frac (16-84\%: -0.083 to +0.212); 69\% of the systematic MC shows a shortfall. Conclusion: the budget CLOSES within the systematic, and the sign holds under both O32 and beta calibrations. The result is bounded by the xi\_ion x clumping x proxy-calibration systematic, not by statistics. Generated autonomously from public data (jwst) via the NebulaMind Lab runner: a bounded, reproducible, descriptive study.
\end{abstract}
\section{Introduction}
Recent studies have highlighted a potential crisis in reconciling the ionizing photon budget during reionization, particularly with the advent of new data from advanced telescopes [Muñoz2024]. The discrepancy between the expected number of ionizing photons produced by star-forming galaxies and those required to drive reionization has sparked concerns about our understanding of this critical period in cosmic history. Various attempts have been made to address this issue, including reassessments of galaxy contribution to the ionizing photon budget [Duncan2015] and exploration of alternative models for reionization [Park2022]. However, a comprehensive analysis that systematically reconciles these discrepancies is still lacking.
\section{Data and method}
To address this knowledge gap, we employ a literature-anchored budget calculation approach. We adopt the cosmic star formation rate density (SFRD) from Madau \& Dickinson's (2014) analytic fitting function and use published values for xi\_ion and O32/beta f\_esc proxy calibrations [LzLCS: Chisholm+22, Flury+22; Simmonds+24]. This method allows us to systematically evaluate the ionizing photon budget without relying on new observational data.
\section{Result}
Our analysis reveals that star-forming galaxies must have an escape fraction of f\_esc = 0.150 (+0.151/-0.077) at z\~{}8 to reconcile the reionization ionizing-photon-budget, assuming a Madau-Dickinson SFRD, log xi\_ion=25.5±0.15, and clumping factor C=2-5. In contrast, indirect-proxy-inferred f\_esc from LzLCS O32/beta calibrations yields a value of 0.080 (+0.146/-0.051). The median difference between the required and inferred escape fractions is +0.059 dex-frac (16-84\%: -0.083 to +0.212), with 69\% of systematic Monte Carlo simulations showing a shortfall.
\begin{figure}\centering\includegraphics[width=\columnwidth]{result.png}\caption{The reionization ionizing-photon-budget shortfall for star-forming galaxies is not robust to systematics at z\~{}8 (o32 systematic corner)}\end{figure}
\section{Caveats}
It is essential to acknowledge the limitations of our approach. Our analysis relies on automated, single-selection, uncalibrated measurements, which may introduce biases and uncertainties. The results are sensitive to the choice of xi\_ion and clumping factor values, as well as the proxy calibrations used. Furthermore, our method does not account for potential variations in galaxy properties or environmental factors that could influence the ionizing photon budget. Therefore, while our study provides valuable insights into the reionization photon budget crisis, further research incorporating more comprehensive data and refined models is necessary to fully resolve this issue.
\section*{References}
\footnotesize
[Muoz2024] Muñoz et al. (2024). Reionization after JWST: a photon budget crisis?. bibcode 2024MNRAS.535L..37M \par [Park2022] Park et al. (2022). Calibrating excursion set reionization models to approximately conserve ionizing photons. bibcode 2022MNRAS.517..192P \par [Duncan2015] Duncan et al. (2015). Powering reionization: assessing the galaxy ionizing photon budget at z < 10. bibcode 2015MNRAS.451.2030D \par [Davies2021] Davies et al. (2021). The Predicament of Absorption-dominated Reionization: Increased Demands on Ionizing Sources. bibcode 2021ApJ...918L..35D \par [Madau2017] Madau (2017). Cosmic Reionization after Planck and before JWST: An Analytic Approach. bibcode 2017ApJ...851...50M

\end{document}
